% Converted on 2026-10-08 from archive/2026-10-08-development-beamer/sections/02-system-workflow.tex (Beamer, review copy 43).
% Each frame became a subsection whose body, unchanged, sits in a group as in Beamer.

\section{시스템 구조와 계산 흐름}



\subsection{모델과 계산 입력 준비}\label{workflow}
\begingroup% frame-local settings, as in Beamer
  \framesubtitle{목표 구조 1/2 · 해밀토니안·계산격자·스핀 배열의 역할을 구분한다.}
  \centering
  \begin{tikzpicture}
    \node[box,font=\normalsize,text width=3.1cm,minimum height=1.3cm] (builders) at (0,1.3)
      {모델별 정의\\NBCP 등};
    \node[focus,font=\normalsize,text width=3.1cm,minimum height=1.3cm] (model) at (4.6,1.3)
      {공통 물리 모델\\해밀토니안};
    \node[focus,font=\normalsize,text width=3.1cm,minimum height=1.3cm] (system) at (9.2,1.3)
      {계산 대상 스핀계\\사이트·결합·계수};
    \node[box,font=\normalsize,text width=3.1cm,minimum height=1.3cm] (geometry) at (9.2,-1.1)
      {계산격자 조건\\크기·형상·경계조건};
    \node[box,font=\normalsize,text width=3.5cm,minimum height=1.3cm] (state) at (4.6,-1.1)
      {스핀 배열\\스핀 방향\\자기구조의 반복주기};
    \draw[flow] (builders)--(model);
    \draw[flow] (model)--(system);
    \draw[flow] (geometry.north)--(system.south);
    \draw[auxflow,{Latex[length=2mm]}-{Latex[length=2mm]}] (state.east)--(geometry.west)
      node[midway,above,font=\small,text=KFBMuted]{정합};
  \end{tikzpicture}
  \par\bigskip
  \begin{minipage}{0.98\textwidth}
    \supportingtext{모델은 해밀토니안이며, 계산 대상 스핀계는 여기에 계산격자 조건을 적용해 구성한다.\\
    스핀 배열은 고전 최적화의 초기 배열 또는 LSWT의 기준 배열로 입력하며, 계산격자와 정합해야 한다.}
  \end{minipage}
  \slidesource{합의된 블록 구조; 공통 모델 규약과 전체 연결은 개발 대상}{contract}
\endgroup


\subsection{계산 방법과 결과 활용}\label{workflow-calculation}
\begingroup% frame-local settings, as in Beamer
  \framesubtitle{목표 구조 2/2 · 계산 대상 스핀계에서 방법을 선택하고 결과를 분석·표현한다.}
  \centering
  \begin{tikzpicture}
    \node[font=\normalsize,text=KFBNavy] (system) at (4.6,3.55)
      {계산 대상 스핀계 · 외부 변수(장 $B$, 온도 $T$)};
    \node[box,font=\normalsize,text width=3.1cm,minimum height=1.3cm] (classical) at (0,1.75)
      {고전 계산\\에너지·최적화};
    \node[box,font=\normalsize,text width=3.1cm,minimum height=1.3cm] (lswt) at (4.6,1.75)
      {LSWT};
    \node[box,dashed,font=\normalsize,text width=3.1cm,minimum height=1.3cm] (quantum) at (9.2,1.75)
      {ED\\스핀파 결과 검증용};
    \node[focus,font=\normalsize,text width=3.1cm,minimum height=1.1cm] (result) at (0,-0.6)
      {계산 결과};
    \node[box,font=\normalsize,text width=3.1cm,minimum height=1.1cm] (observables) at (4.6,-0.6)
      {물리량 분석};
    \node[box,font=\normalsize,text width=3.1cm,minimum height=1.1cm] (output) at (9.2,-0.6)
      {시각화·결과 저장};
    \draw[flow] (system.south)--(lswt.north);
    \draw[flow] (4.6,2.9)--(0,2.9)--(classical.north);
    \draw[flow] (4.6,2.9)--(9.2,2.9)--(quantum.north);
    \fill[KFBBlue] (4.6,2.9) circle (1.25pt);
    \draw[flow] ([yshift=0.25cm]classical.east)--([yshift=0.25cm]lswt.west);
    \node[font=\footnotesize,text=KFBMuted,align=center,anchor=south,inner sep=1pt] at (2.3,2.05)
      {기준\\상태};
    \draw[auxflow] ([yshift=-0.25cm]lswt.west)--([yshift=-0.25cm]classical.east);
    \node[font=\footnotesize,text=KFBMuted,align=center,anchor=north,inner sep=1pt] at (2.3,1.45)
      {영점\\에너지};
    \draw[flow] (classical.south)--(result.north);
    \draw[draw=KFBBlue,line width=0.85pt]
      (lswt.south)--(4.6,0.45)--(0,0.45);
    \draw[draw=KFBBlue,line width=0.85pt]
      (quantum.south)--(9.2,0.45)--(4.6,0.45);
    \fill[KFBBlue] (0,0.45) circle (1.25pt);
    \fill[KFBBlue] (4.6,0.45) circle (1.25pt);
    \draw[flow] (result)--(observables);
    \draw[flow] (observables)--(output);
  \end{tikzpicture}
  \par\bigskip
  \begin{minipage}{0.98\textwidth}
    \supportingtext{범위는 반고전 방법이다(2026-10-10 결정): 고전 탐색, LSWT, $1/S$ 보정, 고전 MC·스핀 동역학. ED는 스핀파 결과 검증용으로만 두고, 텐서 네트워크·DMRG·QMC는 다루지 않는다.\\
    영점 에너지로 고전 manifold 위의 상태를 고르는 경로(D17 궤도 선택)가 있어 LSWT가 상태 선택에 되먹임된다.\\
    LSWT의 기준 상태는 직접 지정할 수도 있다. ED는 고전 최적화를 거치지 않는다.}
  \end{minipage}
  \slidesource{합의된 계산·분석 흐름; 공통 규약을 통한 연결은 후속 개발 대상}{contract}
\endgroup


\subsection{현재 구현의 계산 흐름 (2026-10-21)}\label{workflow-implementation}
\begingroup% frame-local settings, as in Beamer
  \framesubtitle{위 목표 구조가 코드에서 어떤 이름으로 이어지는지. 옛 진입점은 0.2에서 경고, 0.3에서 지운다(D43, D65).}
  \footnotesize\renewcommand{\arraystretch}{1.08}
  \begin{tabularx}{\textwidth}{@{}p{2.6cm}p{4.6cm}Y@{}}
    \toprule
    블록 & 코드 & 규칙 \\
    \midrule
    모델 & \texttt{SpinModel}, \texttt{Site}, \texttt{Term.bilinear}·\texttt{zeeman}·\texttt{onsite} & 만들 때 검증. 값은 $E_0$ 단위(D20). 방법 전용 로직을 넣지 않는다 \\
    외부 조건·설정 & \texttt{ExternalConditions}(장, 온도), \texttt{LSWTSettings}(메시, 정규화, 허용오차) & 조건과 수치 설정을 나눈다 \\
    상태 & \texttt{SpinState}, \texttt{classical\_search}, \texttt{find\_stationary\_state}, \texttt{IncommensurateStructure} & 고전 탐색은 초격자 LT 하한까지 다시 돈다(D70). 단위 세포를 고정한 정상점 탐색(D63) \\
    안정성·상 비교 & \texttt{local\_stability}, \texttt{compare\_states} & 연속 $k$ 국소 안정성(D51), 안정성을 함께 적는 상 비교(D53) \\
    LSWT & \texttt{solve\_lswt} $\to$ \texttt{LSWTResult}, \texttt{solve\_spiral\_lswt} & 불안정이면 \texttt{LSWTError}(D61). $H(k)$는 결합별 희소 저장(D67). 결과는 읽기 전용, 머리에 출처·설정·진단(D14, D66) \\
    선형 너머 & \texttt{solve\_nlswt}, \texttt{pseudo\_goldstone\_gap}, \texttt{solve\_scswt} & $1/S$ 보정(D46, D47). 메시가 풀지 못하는 값은 NaN(D69, D71). 자기무모순 one-loop는 $H_3=0$인 상태만(D72) \\
    물리량 & \texttt{spintoolkit.observables}의 띠, 구조 인자, 중성자 세기, 열역학, 위상량, \texttt{magnetization\_curve} & 발산은 부호 있는 $\pm\infty$, 정의되지 않음은 NaN, 영모드 보손 수는 탐지해 알린다(D64) \\
    검증 & \texttt{spintoolkit.methods.ed.solve\_ed} & 스핀파 결과 검증용 \\
    \bottomrule
  \end{tabularx}
  \slidesource{2026-10-21 코드와 결정 D43--D72}{contract,implementation}
\endgroup


\subsection{구성요소별 책임과 전달 정보}\label{boundaries}
\begingroup% frame-local settings, as in Beamer
  \framesubtitle{기능별 역할을 정하고, 다음 블록이 사용할 정보를 명확하게 전달한다.}
  \begin{tabularx}{\textwidth}{@{}p{3.0cm}Y@{}}
    \toprule
    구성요소 & 담당 기능과 전달 정보 \\
    \midrule
    \textbf{모델 정의} & 해밀토니안: 격자·스핀 크기·항·g-텐서를 공통 형식으로 제공한다. \\
    \textbf{계산계 구성} & 모델에 계산격자 조건을 적용하여 계산 대상 스핀계를 구성한다. \\
    \textbf{상태 정의·최적화} & 스핀 방향·반복주기를 표현하고, 계산계와 정합한 기준 상태를 탐색한다. \\
    \textbf{계산 방법} & 계산계·상태·외부 변수와 방법별 수치 설정을 받아 계산을 수행한다. \\
    \textbf{결과·물리량 분석} & 방법별 결과를 보존하고, 필요한 정보로 물리량을 계산한다. \\
    \textbf{시각화·저장} & 결과를 표현하고 입력 조건·외부 변수·단위와 함께 보관한다. \\
    \bottomrule
  \end{tabularx}
  \par\medskip
  \begin{band}{전달 경계}
    스핀 크기 $S$와 g-텐서는 모델에, 스핀 방향과 자기구조의 반복주기는 상태에 속한다.\\
    장 $B$·온도 $T$는 외부 변수이며, k점 격자·허용오차 등의 수치 설정과 구분한다.
  \end{band}
  \slidesource{합의된 기능별 책임; 외부 변수와 상태 조건은 다음 두 쪽에서 설명}{contract}
\endgroup


\subsection{외부 변수는 해밀토니안과의 관계로 구분한다}\label{external-variables}
\begingroup% frame-local settings, as in Beamer
  \framesubtitle{모델 계수를 바꾸면 새 모델을 만들고, 장·온도를 바꾸면 계산 요청만 바뀐다.}
  \small
  \begin{tabularx}{\textwidth}{@{}p{2.7cm}p{3.9cm}Yp{2.9cm}@{}}
    \toprule
    부류 & 예 & $H$에 들어가는가 & 변경 시 처리 \\
    \midrule
    해밀토니안 파라미터 & 모델 계수 $J$, $K$, $\Gamma$ & 예 & 새 모델 생성 \\
     & 외부 장 $B$ & 예: 모델의 $g$와 결합 & 계산 요청에서 변경 \\
    통계 조건 & 온도 $T$ & 아니오: $\rho\propto e^{-H/k_BT}$ & 계산 요청에서 변경 \\
    계산계 실현 & 열역학 극한; 유한 클러스터 크기·형상·경계조건 & 유한 연산자의 공간과 항 집합을 결정 & 계산계 재구성 \\
    수치 설정 & k점 격자, 허용오차, 최적화 초기값 & 아니오 & 방법 설정에서 변경 \\
    \bottomrule
  \end{tabularx}
  \par\medskip
  \begin{mdcallout}[note]{k점 격자의 두 가지 의미}
    LSWT의 $N_1\times N_2$ k점 격자는 같은 수의 자기 단위격자로 된 주기 경계 토러스와 같다.
    열역학 극한에서는 수렴을 조절하는 수치 설정이지만, 같은 클러스터의 ED와 비교할 때는
    계산계 실현을 정한다. 결과에 어느 의미로 사용했는지 기록한다.
  \end{mdcallout}
  \slidesource{외부 변수 분류; 2026-09-29 검토}{contract}
\endgroup


\subsection{상태는 계산계와 정합해야 한다}\label{state-geometry}
\begingroup% frame-local settings, as in Beamer
  \framesubtitle{스핀 배열은 독립 입력이 아니며, LSWT 기준 상태에는 정상성 조건이 있다.}
  \begin{columns}[T,onlytextwidth]
    \begin{column}{0.48\textwidth}
      \begin{block}{계산계와의 정합}
        \begin{itemize}
          \item 유한 주기 클러스터: 자기 초격자가 클러스터를 빈틈없이 채워야 한다.
            허용되는 질서 파수벡터는 클러스터의 이산 k 집합으로 제한된다.
          \item 예: 120° 구조는 $L$이 3의 배수인 클러스터에만 놓인다.
          \item 열역학 극한 LSWT: 상태의 자기 단위격자가 계산 셀을 정한다.
        \end{itemize}
      \end{block}
      \smallskip
      한쪽 방향의 입력이 아니라 \textbf{정합성 제약}으로 검사한다.
    \end{column}
    \begin{column}{0.48\textwidth}
      \begin{block}{LSWT 기준 상태의 조건}
        \begin{itemize}
          \item 각 스핀이 국소장과 평행(토크 0)해야 보손 선형항이 사라진다.
            극소 여부는 Colpa 조건으로 판정한다.
          \item 정상성은 상대 토크로 판정해 결과 진단에 적고, 정상이 아닌 상태의 결과는
            물리량 함수가 \textbf{거부}한다(D68). 정규화 이동은 불안정한 상태에 쓰지 않는다(D61).
          \item 고전적으로 불안정한 정상 상태는 켤 때만 쓰는 최소 증강 스핀파(MAGSWT,
            \texttt{augmentation="minimal"})로 에너지만 낸다(D62).
          \item 영점 에너지로 상태를 고를 때는 고전 바닥상태 manifold 위에서만 비교하고(D17),
            선택에 쓴 에너지를 결과에 기록한다.
        \end{itemize}
      \end{block}
    \end{column}
  \end{columns}
  \par\medskip
  \supportingtext{ED·TN은 고전 기준 상태를 필수로 요구하지 않는다.}
  \slidesource{상태와 계산계의 관계; 기존 최적화 경로 확인}{contract,implementation}
\endgroup

